\paragraph{Robustness sweeps.}\label{sec:sweeps}
Table~\ref{tab:sweeps} varies one factor at a time on Fashion-MNIST: the strength of the concept shift, the number of groups, the fraction of participating clients and the number of clients. Baselines are given the true $K$ where they need it.

\begin{table}[!htbp]
\centering
\caption{Robustness sweeps (Fashion-MNIST, 3 seeds; the $90^\circ$ row repeats the \emph{Rotation} scenario with 5 seeds). For DisCo-CFL: ARI, number of clusters, conflict-merge rate, accuracy and concept accuracy. For comparison: the best baseline on accuracy and on concept accuracy (in parentheses), and the Oracle's accuracy.}\label{tab:sweeps}
\scriptsize
\setlength{\tabcolsep}{3pt}
\resizebox{\textwidth}{!}{%
\begin{tabular}{llcccccccc}
\toprule
& & \multicolumn{5}{c}{DisCo-CFL} & & & \\
\cmidrule(lr){3-7}
Factor & Setting & ARI & $K$ & Conflict\% & Acc & Concept & Best baseline Acc & Best baseline Concept & Oracle Acc \\
\midrule
Strength & $\theta=15^\circ$ & 0.52 & 2.7 & 30.0 & 88.7 & 76.0 & 88.7 (PACFL) & 75.1 (PACFL) & 89.3 \\
Strength & $\theta=30^\circ$ & 0.93 & 4.0 & 1.7 & 88.8 & 76.3 & 88.3 (FedDAG-C) & 74.5 (PACFL) & 89.1 \\
Strength & $\theta=45^\circ$ & 1.00 & 4.0 & 0.0 & 89.1 & 77.0 & 89.3 (FedDAG-C) & 76.6 (FedDAG-C) & 89.2 \\
Strength & $\theta=90^\circ$ & 0.92 & 4.4 & 1.1 & 89.1 & 76.4 & 88.6 (FedDAG-C) & 74.3 (FL+HC) & 89.3 \\
Groups & $K=2$ & 1.00 & 2.0 & 0.0 & 89.5 & 78.2 & 89.3 (FL+HC) & 77.3 (FL+HC) & 89.5 \\
Groups & $K=4$ & 0.97 & 4.0 & 0.0 & 89.2 & 76.7 & 88.3 (FL+HC) & 72.6 (FL+HC) & 89.2 \\
Groups & $K=8$ & 0.79 & 8.7 & 1.7 & 88.2 & 70.7 & 87.2 (MTCFL) & 59.8 (MTCFL) & 88.6 \\
Participation 30\% & rotation & 0.93 & 4.3 & 0.8 & 87.1 & 72.3 & 86.9 (FedDAG-C) & 70.0 (FedDAG-C) & 87.2 \\
Participation 30\% & label swap & 0.98 & 4.0 & 0.0 & 87.3 & 73.1 & 85.7 (FedCCFA) & 66.4 (FL+HC) & 87.2 \\
Scale & $N=200$ & 0.88 & 4.7 & 2.9 & 86.5 & 73.1 & 85.2 (IFCA) & 70.6 (FedDAG-C) & 87.2 \\
\bottomrule
\end{tabular}}

\end{table}

\paragraph{Strength of the shift.} With rotations of $0,\theta,2\theta,3\theta$ degrees, DisCo-CFL separates the groups perfectly at $\theta=45^\circ$ (ARI $1.00$) and almost perfectly at $30^\circ$ (ARI $0.93$). At $15^\circ$ neighbouring concepts are so close that their class-conditional gradients agree above the threshold: the estimated separation is $\hat\Delta=0.06$, below $1-\tau=0.1$ (Table~\ref{tab:diag}), and DisCo-CFL merges them (ARI $0.52$, $K=2.7$). This is what the separation assumption predicts: the method groups what it cannot tell apart, and the diagnostic shows it. It costs little, because nearby concepts can share a model: accuracy is 88.7\% against 89.3\% for the Oracle, equal to the best baseline (PACFL, given $K$). For $\theta\ge30^\circ$ DisCo-CFL is within $0.3$ points of the Oracle; at $45^\circ$ FedDAG-C is $0.2$ points more accurate.

\paragraph{Number of groups.} With $K=2$, $4$ and $8$ label-permutation groups, DisCo-CFL finds the number of groups without being told (2.0, 4.0 and 8.7 clusters). With $K=8$, recovery is incomplete (ARI $0.79$, 1.7\% of the conflicting pairs merged): with Dirichlet label skew, many clients hold few of the classes that distinguish two of the eight permutations, so coverage is weaker. Accuracy is within $0.4$ points of the Oracle for all $K$, and $0.2$--$1.0$ points above the best baseline; concept accuracy is $0.9$--$10.9$ points above the best baseline.

\paragraph{Partial participation.} When only 30\% of the clients take part in each round, the warm-up model is noisier. DisCo-CFL still recovers the groups (ARI $0.93$ and $0.98$) and matches the Oracle within $0.1$ points, $0.2$--$1.6$ points above the best baseline in accuracy and $2.3$--$6.7$ points in concept accuracy.

\paragraph{Number of clients.} With 200 clients of 200 samples each, DisCo-CFL reaches ARI $0.88$ with 2.9\% of the conflicting pairs merged. Accuracy is 86.5\% against 87.2\% for the Oracle and 85.2\% for the best baseline (IFCA, given $K$); concept accuracy is $2.5$ points above the best baseline.
